
% =========================
% References
% =========================
% --------------- BIBLIOGRAPHY ---------------
\newpage
\addcontentsline{toc}{chapter}{References}
\begin{thebibliography}{99}

% % Thai reference
% \bibitem{book} 
% \thai{นพดล เรียบเลิศหิรัญ}. (2548). 
% \thai{การปลูกพืชไร้ดิน ครั้งที่พิมพ์. สำนักพิมพ์ กรุงเทพฯ}, หน้า 10–15.


% % English journal example
% \bibitem{journal} 
% Nooijen, W.F.J.M. and Muilwijk, B. (2005). 
% "Paint/Water Separation by Ceramic Microfiltration". 
% \textit{Environmental Progress}, 31(2), 86–96.


% Hinton (2025)
\bibitem{fortune2025} 
Hinton, G. (2025). “Give AI a maternal instinct so it cares for us.”  
\textit{Fortune Magazine}, August 2025.  
Retrieved from \url{https://fortune.com/2025/08/14/godfather-of-ai-geoffrey-hinton-maternal-instincts-superintelligence/}.

% Hamilton A
\bibitem{hamilton1964a}
Hamilton, W. D. (1964).
“The genetical evolution of social behaviour. I.”
\textit{Journal of Theoretical Biology}, 7(1), 1–16.
Retrieved from \url{https://doi.org/10.1016/0022-5193(64)90038-4}.
% Hamilton B
\bibitem{hamilton1964b}
Hamilton, W. D. (1964).
“The genetical evolution of social behaviour. II.”
\textit{Journal of Theoretical Biology}, 7(1), 17–52.
Retrieved from \url{https://doi.org/10.1016/0022-5193(64)90039-6}.

% Gorrell2010
\bibitem{gorrell2010}
Gorrell, J. C., McAdam, A. G., Coltman, D. W., Humphries, M. M., \& Boutin, S. (2010).
``Adoption of kin enhances inclusive fitness in asocial red squirrels.''
\textit{Nature Communications}, 1, 22.

\bibitem{kol liker2015}
Kölliker, M., Royle, N. J., Smiseth, P. T., et al. (2015).
``Parent–offspring conflict and the evolution of family life.''
\textit{Nature Communications}, 6, 7220.


\bibitem{haig2014}
Haig, D. (2014).
``Coadaptation and conflict, misconception and muddle, in the evolution of genomic imprinting.''
\textit{Heredity}, 113, 96–103.

\bibitem{cluttonbrock2002}
Clutton-Brock, T. H. (2002).
``Breeding together: kin selection and mutualism in cooperative vertebrates.''
\textit{Science}, 296(5565), 69–72.

\bibitem{cockburn2006}
Cockburn, A. (2006).
``Prevalence of different modes of parental care in birds.''
\textit{Proceedings of the Royal Society B: Biological Sciences}, 273, 1375–1383.

\bibitem{griffe2003}
Griffin, A. S., \& West, S. A. (2003).
``Kin discrimination and the benefit of helping in cooperatively breeding vertebrates.''
\textit{Science}, 302(5645), 634–636.

\bibitem{hodge2007}
Hodge, S. J. (2007).
``Associations between helper number and reproductive success in a cooperative breeder: testing the assumptions of the ‘load-lightening’ hypothesis.''
\textit{Behavioral Ecology and Sociobiology}, 62, 913–922.

\bibitem{komdeur1992}
Komdeur, J. (1992).
``Importance of habitat saturation and territory quality for evolution of cooperative breeding in the Seychelles warbler.''
\textit{Nature}, 358, 493–495.

\bibitem{hrdy1979}
Hrdy, S. B. (1979).
``Infanticide among animals: a review, classification, and examination of the implications for the reproductive strategies of females.''
\textit{Ethology and Sociobiology}, 1(1), 13–40.

\bibitem{stockley2011}
Stockley, P., \& Bro‐Jørgensen, J. (2011).
``Female competition and its evolutionary consequences in mammals.''
\textit{Biological Reviews}, 86(2), 341–366.

\bibitem{lukas2014}
Lukas, D., \& Huchard, E. (2014).
``The evolution of infanticide by females in mammals.''
\textit{Nature Communications}, 5, 4972.


\bibitem{cluttonbrock2007}
Clutton-Brock, T. H. (2007).
``Sexual selection in males and females.'' 
\textit{Science}, 318(5858), 1882–1885.

\bibitem{sapolsky2005}
Sapolsky, R. M. (2005).
``The influence of social hierarchy on primate health.''
\textit{Science}, 308(5722), 648–652.

\bibitem{creel2001}
Creel, S. (2001).
``Social dominance and stress hormones.''
\textit{Trends in Ecology \& Evolution}, 16(9), 491–497.

\bibitem{johnstone2000}
Johnstone, R. A. (2000).
``Models of reproductive skew: a review and synthesis.''
\textit{Trends in Ecology \& Evolution}, 15(7), 418–425.


\bibitem{emlen1977}
Emlen, S. T., \& Oring, L. W. (1977).
``Ecology, sexual selection, and the evolution of mating systems.''
\textit{Science}, 197(4300), 215–223.

\bibitem{lukas2018}
Lukas, D., \& Clutton-Brock, T. H. (2018).
``Social complexity and kinship in animal societies.''
\textit{Ecology Letters}, 21(8), 1129–1134.


% Competitive: Model
\bibitem{maynardsmith1973}
Maynard Smith, J., \& Price, G. R. (1973).
``The logic of animal conflict.''
\textit{Nature}, 246, 15–18.

\bibitem{nowak2004}
Nowak, M. A., \& Sigmund, K. (2004).
``Evolutionary dynamics of biological games.''
\textit{Science}, 303(5659), 793–799.

\bibitem{leibo2017}
Leibo, J. Z., Zambaldi, V., Lanctot, M., Marecki, J., \& Graepel, T. (2017).
``Multi-agent reinforcement learning in sequential social dilemmas.''
\textit{Proceedings of the 16th International Conference on Autonomous Agents and Multiagent Systems (AAMAS)}, 464–473.

\bibitem{wang2018}
Wang, J. X., Hughes, E., Fernando, C., Czarnecki, W. M., Duéñez-Guzmán, E. A., \& Leibo, J. Z. (2019).
``Evolving intrinsic motivations for altruistic behavior.''
\textit{Proceedings of the 18th International Conference on Autonomous Agents and Multiagent Systems (AAMAS)}.
Preprint: \url{https://arxiv.org/abs/1811.05931}.


% Biological mechanisms
\bibitem{dulac2014}
Dulac, C., O'Connell, L. A., \& Wu, Z. (2014).
``Neural control of maternal and paternal behaviors.''
\textit{Science}, 345(6198), 765–770.

\bibitem{galef2013}
Galef, B. G., \& Laland, K. N. (2013).
``Social learning in animals: Empirical studies and theoretical models.''
\textit{BioScience}, 63(6), 465–475.

\bibitem{champagne2006}
Champagne, F. A., \& Meaney, M. J. (2006).
``Transgenerational effects of social environment on variations in maternal care and behavioral response to novelty.''
\textit{Behavioral Neuroscience}, 120(5), 1355–1363.

\bibitem{meaney2010}
Meaney, M. J. (2010).
``Epigenetics and the biological definition of gene × environment interactions.''
\textit{Child Development}, 81(1), 41–79.

\bibitem{numan2006}
Numan, M., \& Insel, T. R. (2006).
\textit{The Neurobiology of Parental Behavior.}
Springer, New York.

\bibitem{bales2015}
Bales, K. L., \& Saltzman, W. (2015).
``Parenting as a model: Neuroendocrine regulation of the transition to caregiving.''
\textit{Hormones and Behavior}, 77, 98–112.

\bibitem{rosenblatt1994}
Rosenblatt, J. S. (1994).
``Psychobiology of maternal behavior: Contributions from the rat.''
\textit{American Journal of Physiology}, 266, R875–R889.

\bibitem{numan2018}
Numan, M. (2018).
``The parental brain: Mechanisms, development, and evolution.''
\textit{Current Topics in Behavioral Neurosciences}, 43, 1–24.


% Env & social factor
\bibitem{cluttonbrock2016}
Clutton-Brock, T. H. (2016).
\textit{Mammal Societies.}
Wiley-Blackwell.

\bibitem{ebensperger2001}
Ebensperger, L. A. (2001).
``A review of the evolutionary causes of rodent group-living.''
\textit{Acta Theriologica}, 46(2), 115–144.

\bibitem{ridley2007}
Ridley, A. R., \& Raihani, N. J. (2007).
``Variable postfledging care in a cooperative bird: Causes and consequences.''
\textit{Behavioral Ecology}, 18(6), 994–1000.

\bibitem{mcnamara2008}
McNamara, J. M., \& Houston, A. I. (2008).
``Optimal annual routines: Behaviour in the context of physiology and ecology.''
\textit{Philosophical Transactions of the Royal Society B: Biological Sciences}, 363(1490), 301–319.

% Evolution theory of altruism 
\bibitem{trivers1971}
Trivers, R. L. (1971).
``The evolution of reciprocal altruism.''
\textit{Quarterly Review of Biology}, 46(1), 35–57.

\bibitem{axelrod1981}
Axelrod, R., \& Hamilton, W. D. (1981).
``The evolution of cooperation.''
\textit{Science}, 211(4489), 1390–1396.

\bibitem{carter2013}
Carter, G. G., \& Wilkinson, G. S. (2013).
``Food sharing in vampire bats: reciprocal help predicts donations more than relatedness or harassment.''
\textit{Proceedings of the Royal Society B: Biological Sciences}, 280(1753), 20122573.

\bibitem{bshary2008}
Bshary, R., \& Grutter, A. S. (2008).
``Asymmetric cheating opportunities and partner control in a cleaner fish mutualism.''
\textit{Animal Behaviour}, 76(2), 439–445.

\bibitem{wilson2007}
Wilson, D. S., \& Wilson, E. O. (2007).
``Rethinking the theoretical foundation of sociobiology.''
\textit{Quarterly Review of Biology}, 82(4), 327–348.

\bibitem{okasha2006}
Okasha, S. (2006).
\textit{Evolution and the Levels of Selection.}
Oxford University Press, Oxford.

\bibitem{nowak2006}
Nowak, M. A. (2006).
``Five rules for the evolution of cooperation.''
\textit{Science}, 314(5805), 1560–1563.

% Computational Model
\bibitem{wilensky1999}
Wilensky, U. (1999).
\textit{NetLogo.} Center for Connected Learning and Computer-Based Modeling, Northwestern University.

\bibitem{west2011}
West, S. A., Griffin, A. S., \& Gardner, A. (2011).
``Evolutionary explanations for cooperation.''
\textit{Current Biology}, 21(19), R959–R971.

\bibitem{hughes2018}
Hughes, C., Leibo, J. Z., Phillips, M., et al. (2018).
``Inequity aversion improves cooperation in intertemporal social dilemmas.''
\textit{Advances in Neural Information Processing Systems (NeurIPS)}, 31.

\bibitem{jaques2019}
Jaques, N., Lazaridou, A., Hughes, E., et al. (2019).
``Social influence as intrinsic motivation for multi-agent deep reinforcement learning.''
\textit{Proceedings of the 36th International Conference on Machine Learning (ICML)}, 3040–3049.

\bibitem{breazeal2003}
Breazeal, C. (2003).
``Toward sociable robots.''
\textit{Robotics and Autonomous Systems}, 42(3–4), 167–175.

\bibitem{kuntoro2022}
Kuntoro, W., et al. (2022).
``Emotion-based caregiving behaviour in assistive robotics: A bio-inspired framework.''
\textit{Frontiers in Robotics and AI}, 9, 873625.

\bibitem{trianni2008}
Trianni, V., \& Dorigo, M. (2008).
``Self-organisation and communication in groups of simulated and physical robots.''
\textit{Biological Cybernetics}, 99(3), 183–200.

\bibitem{hamann2018}
Hamann, H. (2018).
\textit{Swarm Robotics: A Formal Approach.}
Springer, Cham.

% Plasticity, learning, homeostasis, and evolutionary optimization
\bibitem{hebb1949}
Hebb, D. O. (1949).
\textit{The Organization of Behavior: A Neuropsychological Theory.}
Wiley, New York.

\bibitem{cannon1932}
Cannon, W. B. (1932).
\textit{The Wisdom of the Body.}
W. W. Norton \& Company, New York.

\bibitem{holland1975}
Holland, J. H. (1975).
\textit{Adaptation in Natural and Artificial Systems.}
University of Michigan Press, Ann Arbor.

\bibitem{goldberg1989}
Goldberg, D. E. (1989).
\textit{Genetic Algorithms in Search, Optimization, and Machine Learning.}
Addison-Wesley, Reading, MA.

\bibitem{rechenberg1973}
Rechenberg, I. (1973).
\textit{Evolutionsstrategie: Optimierung technischer Systeme nach Prinzipien der biologischen Evolution.}
Frommann-Holzboog, Stuttgart.

\bibitem{schwefel1977}
Schwefel, H.-P. (1977).
\textit{Numerische Optimierung von Computer-Modellen mittels der Evolutionsstrategie.}
Birkhäuser, Basel.

% Maroto-Gómez et al. (2024)
\bibitem{marotogomez2024}
Maroto-Gómez, M., Bueno-Adrada, M., Malfaz, M., Castro-González, Á., Salichs, M. Á. (2024).
“Human–robot pair-bonding from a neuroendocrine perspective: Modeling the effect of oxytocin, arginine vasopressin, and dopamine on the social behavior of an autonomous robot.”
\textit{Robotics and Autonomous Systems}, 176, 104687.
Retrieved from \url{https://doi.org/10.1016/j.robot.2024.104687}.

\end{thebibliography}


% =========================
% Biography
% =========================
\newpage
\chapter*{Biography}
\addcontentsline{toc}{chapter}{Biography}

\noindent\textbf{Name:} Mr.\ Napat Aeimwiratchai \\[0.3cm]
\noindent\textbf{Date of Birth:} 29 March 2004 \\[0.3cm]
\noindent\textbf{Place of Birth:} Thailand \\[0.3cm]
\noindent\textbf{Education:} \\
\begin{itemize}
    \item Bachelor of Engineering (Robotics and Automation Engineering),\\
          Institute of Field Robotics,\\
          King Mongkut's University of Technology Thonburi, 2025
\end{itemize}

\vspace{0.5cm}
\noindent\textbf{Contact Address:} Institute of Field Robotics, King Mongkut's University of Technology Thonburi, 126 Pracha Uthit Road, Bang Mod, Thung Khru, Bangkok 10140, Thailand\\
